\chapter{The Harmonic Oscillator}
\label{ch:sho}

Near the minimum of any potential, $V(x)\approx V(x_0) + \frac12V''(x_0)(x-x_0)^2$, so the harmonic oscillator models every small vibration, from a molecular bond (the step in the heat capacity of diatomic gases) to phonons and photons.

\[
    \hat{H} = \frac{\hat{p}^2}{2m} + \frac12m\omega^2\hat{x}^2 .
\]

\begin{definition}[Natural scales]
Balancing $\hbar^2/2mx_0^2$ against $\frac12m\omega^2x_0^2$ gives
\[
    x_0 = \sqrt{\frac{\hbar}{m\omega}},
    \qquad
    p_0 = \frac{\hbar}{x_0} = \sqrt{\hbar m\omega},
    \qquad
    \hat{H} = \frac{\hbar\omega}{2}\Bigl[\Bigl(\frac{\hat{p}}{p_0}\Bigr)^2 + \Bigl(\frac{\hat{x}}{x_0}\Bigr)^2\Bigr].
\]
\end{definition}

% ============================================================
\section{Solving the Differential Equation}
% ============================================================

With $u = x/x_0$ and $\varepsilon = 2E/\hbar\omega$, the Schr\"odinger equation becomes
\[
    \psi''(u) = (u^2-\varepsilon)\psi(u).
\]
For $u^2\gg\varepsilon$ (deep in the forbidden region), $\psi''\approx u^2\psi$, solved by $e^{\pm u^2/2}$; normalizability picks $e^{-u^2/2}$. Write $\psi = h(u)e^{-u^2/2}$, so that
\[
    h'' - 2uh' + (\varepsilon-1)h = 0 .
\]
A power series $h = \sum_ja_ju^j$ gives the recursion (substituting, $h'' = \sum_j(j+2)(j+1)a_{j+2}u^j$ and $-2uh' = -\sum_j2ja_ju^j$, and requiring each power of $u$ to vanish)
\[
    a_{j+2} = \frac{2j+1-\varepsilon}{(j+1)(j+2)}a_j .
\]
If the series did not terminate, $a_{j+2}\approx\frac2ja_j$ for large $j$, so $h\sim e^{u^2}$ and $\psi\sim e^{u^2/2}$, which is not normalizable. So the series must terminate: $\varepsilon = 2n+1$, with the other parity series set to zero ($a_1 = 0$ for even $n$, $a_0 = 0$ for odd $n$). Then $h$ is a polynomial, the Hermite polynomial $H_n(u) = (-1)^ne^{u^2}\frac{d^n}{du^n}e^{-u^2}$: $H_0 = 1$, $H_1 = 2u$, $H_2 = 4u^2-2$, $H_3 = 8u^3-12u$.

\paragraph{Checking the substitution.} With $\psi = he^{-u^2/2}$, $\psi' = (h'-uh)e^{-u^2/2}$ and $\psi'' = (h'' - 2uh' - h + u^2h)e^{-u^2/2}$. Setting $\psi'' = (u^2-\varepsilon)\psi$, the $u^2h$ terms cancel, leaving $h'' - 2uh' + (\varepsilon-1)h = 0$. For example $n = 2$: with $\varepsilon = 5$, $a_2 = \frac{1-5}{2}a_0 = -2a_0$ and $a_4 = 0$, so $h\propto1-2u^2\propto H_2$.

\begin{formula}[Harmonic oscillator eigenstates]
\[
    E_n = \hbar\omega\bigl(n+\tfrac12\bigr),
    \qquad
    \varphi_n(x) = \frac{1}{\sqrt{2^nn!}}\frac{1}{\pi^{1/4}x_0^{1/2}}H_n\Bigl(\frac{x}{x_0}\Bigr)e^{-x^2/2x_0^2},
    \qquad
    n = 0,1,2,\ldots
\]
\end{formula}

\begin{figure}[ht]
\centering
\begin{tikzpicture}
\begin{axis}[width=9cm,height=7cm,xlabel={$x/x_0$},ylabel={$E/\hbar\omega$},axis lines=left,xmin=-4.2,xmax=4.2,ymin=0,ymax=5,domain=-4.2:4.2,samples=150]
\addplot[thick,black]{x^2/2};
\addplot[thick,blue!70!black]{0.5+0.6*exp(-x^2/2)};
\addplot[thick,blue!70!black]{1.5+0.6*1.41*x*exp(-x^2/2)};
\addplot[thick,blue!70!black]{2.5+0.6*0.71*(2*x^2-1)*exp(-x^2/2)};
\addplot[thick,blue!70!black]{3.5+0.6*0.29*(2*x^3-3*x)*exp(-x^2/2)};
\foreach \e in {0.5,1.5,2.5,3.5} {\edef\t{\noexpand\addplot[dotted,gray]{\e};}\t}
\end{axis}
\end{tikzpicture}
\caption{The harmonic oscillator potential with its first four eigenfunctions, each drawn at its energy $\hbar\omega(n+\frac12)$. The wavefunctions leak into the classically forbidden region beyond the turning points.}
\end{figure}

The levels are equally spaced, so each level above the ground state can be read as $n$ excitations (quanta) of energy $\hbar\omega$ each: phonons for mechanical oscillations, photons for the EM field. The ground state has the \emph{zero-point energy} $\hbar\omega/2\neq0$: $\langle x\rangle = \langle p\rangle = 0$, but $\langle x^2\rangle$ and $\langle p^2\rangle$ cannot both vanish.

\begin{figure}[ht]
\centering
\begin{tikzpicture}
\begin{axis}[width=10cm,height=5cm,xlabel={$x/x_0$},ylabel={probability density},axis lines=left,xmin=-5.5,xmax=5.5,ymin=0,ymax=0.45,samples=400,legend style={draw=none,font=\small,at={(0.5,1)},anchor=north}]
\addplot[thick,blue!70!black,domain=-5.5:5.5]{(1/(2^10*3628800*sqrt(pi)))*(1024*x^10-23040*x^8+161280*x^6-403200*x^4+302400*x^2-30240)^2*exp(-x^2)};
\addlegendentry{quantum, $n=10$}
\addplot[thick,dashed,red!70!black,domain=-4.57:4.57]{1/(pi*sqrt(21-x^2))};
\addlegendentry{classical}
\end{axis}
\end{tikzpicture}
\caption{Position probability density for $n = 10$ compared with the classical density $1/\pi\sqrt{x_{\max}^2-x^2}$ at the same energy. Averaged over the rapid oscillations they agree: the correspondence principle. The classical particle spends most time near the turning points, where it moves slowly.}
\end{figure}

% ============================================================
\section{Ladder Operators}
\label{sec:ladder}
% ============================================================

To factorize $\hat{H}$ as one would $A^2+B^2 = (A-iB)(A+iB)$, we must account for $[\hat{x},\hat{p}] = i\hbar$:
\[
    \Bigl(\frac{\hat{x}}{x_0} - \frac{i\hat{p}}{p_0}\Bigr)\Bigl(\frac{\hat{x}}{x_0} + \frac{i\hat{p}}{p_0}\Bigr) = \frac{\hat{x}^2}{x_0^2} + \frac{\hat{p}^2}{p_0^2} + \frac{i}{x_0p_0}[\hat{x},\hat{p}] = \frac{\hat{x}^2}{x_0^2} + \frac{\hat{p}^2}{p_0^2} - 1 .
\]

\begin{definition}[Ladder operators]
\[
    \hat{a} = \frac{1}{\sqrt2}\Bigl(\frac{\hat{x}}{x_0} + \frac{i\hat{p}}{p_0}\Bigr),
    \qquad
    \hat{a}^\dagger = \frac{1}{\sqrt2}\Bigl(\frac{\hat{x}}{x_0} - \frac{i\hat{p}}{p_0}\Bigr),
    \qquad
    \hat{N} = \hat{a}^\dagger\hat{a}.
\]
Here $\hat{A}^\dagger$ is the adjoint, defined by $\langle f,\hat{A}g\rangle = \langle\hat{A}^\dagger f,g\rangle$; $\hat{x}$ and $\hat{p}$ are their own adjoints, $\partial_x^\dagger = -\partial_x$ and $i^\dagger = -i$. $\hat{a}$ is the annihilation (lowering) and $\hat{a}^\dagger$ the creation (raising) operator.
\end{definition}

\begin{formula}[Oscillator in ladder operators]
\[
    \hat{H} = \hbar\omega\bigl(\hat{N}+\tfrac12\bigr),
    \qquad
    [\hat{a},\hat{a}^\dagger] = 1,
    \qquad
    \hat{x} = \frac{x_0}{\sqrt2}(\hat{a}+\hat{a}^\dagger),
    \qquad
    \hat{p} = \frac{p_0}{\sqrt2\,i}(\hat{a}-\hat{a}^\dagger).
\]
\end{formula}

\begin{fact}[Commutator rules]
$[\hat{A}\hat{B},\hat{C}] = \hat{A}[\hat{B},\hat{C}] + [\hat{A},\hat{C}]\hat{B}$, like a product rule. Hence
\[
    [\hat{N},\hat{a}] = -\hat{a},
    \qquad
    [\hat{N},\hat{a}^\dagger] = \hat{a}^\dagger,
    \qquad
    [\hat{a},(\hat{a}^\dagger)^n] = n(\hat{a}^\dagger)^{n-1}.
\]
\end{fact}

\begin{theorem}[Ladder property]
If $\hat{N}\psi = \lambda\psi$, then $\hat{N}(\hat{a}\psi) = (\lambda-1)\hat{a}\psi$ and $\hat{N}(\hat{a}^\dagger\psi) = (\lambda+1)\hat{a}^\dagger\psi$.
\end{theorem}
\begin{proof}
$\hat{N}\hat{a}\psi = \hat{a}\hat{N}\psi + [\hat{N},\hat{a}]\psi = (\lambda-1)\hat{a}\psi$, and similarly for $\hat{a}^\dagger$.
\end{proof}

The ladder cannot go down forever: for any $\psi$, $\langle\hat{N}\rangle = \langle\psi,\hat{a}^\dagger\hat{a}\psi\rangle = \|\hat{a}\psi\|^2\ge0$, so $\langle\hat{H}\rangle\ge\hbar\omega/2$, with equality only if $\hat{a}\psi = 0$. That condition is a first-order differential equation:
\[
    \Bigl(\frac{x}{x_0} + x_0\frac{d}{dx}\Bigr)\varphi_0 = 0
    \quad\Longrightarrow\quad
    \varphi_0(x) = \frac{1}{\pi^{1/4}x_0^{1/2}}e^{-x^2/2x_0^2}.
\]

\begin{formula}[Oscillator states from the ladder]
\[
    \varphi_n = \frac{(\hat{a}^\dagger)^n}{\sqrt{n!}}\varphi_0,
    \qquad
    \hat{a}^\dagger\varphi_n = \sqrt{n+1}\,\varphi_{n+1},
    \qquad
    \hat{a}\varphi_n = \sqrt{n}\,\varphi_{n-1},
    \qquad
    \hat{N}\varphi_n = n\varphi_n .
\]
\end{formula}
The normalization follows from $C_n = \|(\hat{a}^\dagger)^n\varphi_0\|^2 = \langle\varphi_0,\hat{a}^n(\hat{a}^\dagger)^n\varphi_0\rangle = nC_{n-1}$, using $[\hat{a},(\hat{a}^\dagger)^n] = n(\hat{a}^\dagger)^{n-1}$ and $\hat{a}\varphi_0 = 0$, so $C_n = n!$.

\begin{proof}[The algebra in detail]
Using $\hat{x} = \frac{x_0}{\sqrt2}(\hat{a}+\hat{a}^\dagger)$, $\hat{p} = \frac{p_0}{\sqrt2i}(\hat{a}-\hat{a}^\dagger)$ and $x_0p_0 = \hbar$:
\[
    \frac{\hat{x}^2}{x_0^2} + \frac{\hat{p}^2}{p_0^2} = \frac12\bigl[(\hat{a}+\hat{a}^\dagger)^2 - (\hat{a}-\hat{a}^\dagger)^2\bigr] = \hat{a}\hat{a}^\dagger + \hat{a}^\dagger\hat{a} = 2\hat{a}^\dagger\hat{a} + 1,
\]
using $[\hat{a},\hat{a}^\dagger] = 1$, so $\hat{H} = \frac{\hbar\omega}{2}(2\hat{N}+1)$. The commutator itself is
\[
    [\hat{a},\hat{a}^\dagger] = \frac12\Bigl[\frac{\hat{x}}{x_0} + \frac{i\hat{p}}{p_0},\ \frac{\hat{x}}{x_0} - \frac{i\hat{p}}{p_0}\Bigr] = \frac{1}{2x_0p_0}\bigl(-i[\hat{x},\hat{p}] + i[\hat{p},\hat{x}]\bigr) = \frac{1}{2\hbar}(\hbar+\hbar) = 1 .
\]
Finally $\hat{a}^\dagger\varphi_n = c\,\varphi_{n+1}$ with $|c|^2 = \langle\varphi_n|\hat{a}\hat{a}^\dagger|\varphi_n\rangle = \langle\varphi_n|\hat{N}+1|\varphi_n\rangle = n+1$, and similarly $|\langle\varphi_{n-1}|\hat{a}|\varphi_n\rangle|^2 = n$.
\end{proof}

\begin{example}[Building $\varphi_1$ and $\varphi_2$]
In position space $\hat{a}^\dagger = \frac{1}{\sqrt2}\bigl(\frac{x}{x_0} - x_0\frac{d}{dx}\bigr)$. Acting on $\varphi_0\propto e^{-x^2/2x_0^2}$,
\[
    \hat{a}^\dagger\varphi_0 = \frac{1}{\sqrt2}\Bigl(\frac{x}{x_0} + \frac{x}{x_0}\Bigr)\varphi_0 = \sqrt2\frac{x}{x_0}\varphi_0 = \varphi_1,
\]
and $\hat{a}^\dagger\varphi_1 = \frac{1}{\sqrt2}\bigl(\frac{2x^2}{x_0^2}-1\bigr)\varphi_0\cdot\sqrt2 = \sqrt2\,\varphi_2$, so $\varphi_2\propto(2u^2-1)e^{-u^2/2}\propto H_2(u)e^{-u^2/2}$, matching the series solution.
\end{example}

\begin{example}[Uncertainty in $\varphi_n$]
$\langle\hat{x}\rangle_n = \frac{x_0}{\sqrt2}\langle\varphi_n,(\hat{a}+\hat{a}^\dagger)\varphi_n\rangle = 0$, since $\hat{a}\varphi_n$ and $\hat{a}^\dagger\varphi_n$ are orthogonal to $\varphi_n$. Similarly
\[
    \langle\hat{x}^2\rangle_n = \frac{x_0^2}{2}\|(\hat{a}+\hat{a}^\dagger)\varphi_n\|^2 = \frac{x_0^2}{2}\bigl\|\sqrt{n}\varphi_{n-1} + \sqrt{n+1}\varphi_{n+1}\bigr\|^2 = x_0^2\bigl(n+\tfrac12\bigr),
\]
and likewise $\langle\hat{p}^2\rangle_n = p_0^2(n+\tfrac12)$, so $\Delta x\,\Delta p = \hbar(n+\tfrac12)$: only the ground state saturates the uncertainty bound.
\end{example}

\begin{example}[The half oscillator]
For $V = \frac12m\omega^2x^2$ on $x>0$ and $\infty$ for $x<0$, the wall imposes $\psi(0) = 0$. On $x>0$ the oscillator equation is unchanged, so the solutions are the full oscillator's eigenfunctions that already vanish at the origin: the odd ones. Hence $E = \hbar\omega(n+\frac12)$ for $n = 1,3,5,\ldots$, i.e.\ $\frac32\hbar\omega,\frac72\hbar\omega,\ldots$, with spacing $2\hbar\omega$, and normalized eigenfunctions $\sqrt2\varphi_n$ on the half line.
\end{example}

\begin{example}[Higher moments]
Writing $\hat{x}^2\varphi_n = \frac{x_0^2}{2}\bigl(\sqrt{n(n-1)}\varphi_{n-2} + (2n+1)\varphi_n + \sqrt{(n+1)(n+2)}\varphi_{n+2}\bigr)$ and using orthonormality,
\[
    \langle\hat{x}^4\rangle_n = \|\hat{x}^2\varphi_n\|^2 = \frac{x_0^4}{4}\bigl(6n^2+6n+3\bigr).
\]
The ladder operators reduce every polynomial expectation value to counting.
\end{example}

\begin{example}[A breathing superposition]
For $|\Psi(0)\rangle = \frac{1}{\sqrt2}(|0\rangle + |2\rangle)$, each component picks up its own phase, $|\Psi(t)\rangle = \frac{1}{\sqrt2}\bigl(e^{-i\omega t/2}|0\rangle + e^{-5i\omega t/2}|2\rangle\bigr)$. Since $\hat{x}$ changes $n$ by one, $\langle\hat{x}\rangle = 0$ at all times, but $\hat{x}^2$ connects $|0\rangle$ and $|2\rangle$:
\[
    \langle\hat{x}^2\rangle(t) = \frac{x_0^2}{4}\bigl(6 + 2\sqrt2\cos2\omega t\bigr).
\]
The packet stays centred but its width oscillates at $2\omega$, the Bohr frequency $(E_2-E_0)/\hbar$.
\end{example}

\begin{fact}[$\hat{a}^\dagger$ has no eigenstates]
$\hat{a}$ has eigenstates for every $\alpha\in\mathbb{C}$ (the coherent states of Chapter~\ref{ch:dynamics}). But solving $\hat{a}^\dagger\psi = \alpha\psi$ gives $\psi\propto\exp\bigl(\frac{x^2}{2x_0^2} - \frac{\sqrt2\alpha x}{x_0}\bigr)$, which is not normalizable.
\end{fact}

% ============================================================
\section{Thermal Oscillators and the Three-Dimensional Oscillator}
% ============================================================

At temperature $T$, with $\beta = 1/k_BT$, the partition function $Z = \sum_ne^{-\beta E_n}$ gives $\langle E\rangle = -\partial_\beta\ln Z$ and $C_V = k_B\beta^2\partial_\beta^2\ln Z$. For the oscillator,
\[
    Z = \frac{e^{-\beta\hbar\omega/2}}{1-e^{-\beta\hbar\omega}},
    \qquad
    C_V = k_B\frac{(\beta\hbar\omega)^2e^{-\beta\hbar\omega}}{(1-e^{-\beta\hbar\omega})^2}
    \approx\begin{cases}k_B\,(\beta\hbar\omega)^2e^{-\beta\hbar\omega}, & k_BT\ll\hbar\omega,\\ k_B, & k_BT\gg\hbar\omega.\end{cases}
\]
The vibration is frozen out until $k_BT\sim\hbar\omega$, then contributes the classical $k_B$: the steps in the heat capacity of diatomic gases (Chapter~\ref{ch:experiment}). For H$_2$, estimating $\frac12m\omega^2a_0^2\sim1$\,Ry gives $\hbar\omega/k_B\sim10^4$\,K.

\begin{figure}[ht]
\centering
\begin{tikzpicture}
\begin{axis}[width=9cm,height=5cm,xlabel={$k_BT/\hbar\omega$},ylabel={$C_V/k_B$},axis lines=left,xmin=0,xmax=2.5,ymin=0,ymax=1.1,domain=0.02:2.5,samples=200]
\addplot[thick,blue!70!black]{(1/x)^2*exp(-1/x)/(1-exp(-1/x))^2};
\addplot[dashed,gray,domain=0:2.5]{1};
\end{axis}
\end{tikzpicture}
\caption{Heat capacity of a quantum oscillator. It is exponentially suppressed for $k_BT\ll\hbar\omega$ and approaches the classical equipartition value $k_B$ for $k_BT\gg\hbar\omega$.}
\end{figure}

\begin{example}[Two-dimensional isotropic oscillator]
$\hat{H} = \hat{H}_x + \hat{H}_y$ gives $E = \hbar\omega(n_x+n_y+1) = \hbar\omega(N+1)$, with degeneracy $N+1$. The degeneracy reflects a symmetry larger than rotations: the operators $\hat{a}_i^\dagger\hat{a}_j$ all commute with $\hat{H}$ and move quanta between the $x$ and $y$ oscillators without changing $N$. In terms of circular quanta $\hat{a}_\pm = (\hat{a}_x\mp i\hat{a}_y)/\sqrt2$, the angular momentum is $\hat{L}_z = \hbar(\hat{N}_+ - \hat{N}_-)$, so level $N$ contains $L_z = -N\hbar,-(N-2)\hbar,\ldots,N\hbar$.
\end{example}

\begin{figure}[ht]
\centering
\begin{tikzpicture}[yscale=0.8]
  \foreach \N in {0,1,2,3} {\pgfmathtruncatemacro{\d}{\N+1}\foreach \k in {1,...,\d} {\draw[thick] (\k*0.5-0.2,\N) -- (\k*0.5+0.2,\N);} \node[left] at (0,\N) {\small $N=\N$}; \node[right] at (2.3,\N) {\small $E=\hbar\omega(\N+1)$};}
\end{tikzpicture}
\caption{Levels of the two-dimensional isotropic oscillator, each drawn with its $N+1$ degenerate states.}
\end{figure}

\begin{example}[Isotropic oscillator in three dimensions]
$\hat{H} = \frac{\hat{\lvect{p}}^2}{2m} + \frac12m\omega^2r^2$ separates in Cartesian coordinates, so $E = \hbar\omega(n_x+n_y+n_z+\frac32) = \hbar\omega(N+\frac32)$ with degeneracy $\frac12(N+1)(N+2)$. In spherical coordinates, with $s = r/x_0$ and $u = s^{\ell+1}e^{-s^2/2}w(s)$, the radial equation gives a terminating series for $w$ only if
\[
    E = \hbar\omega\bigl(2k+\ell+\tfrac32\bigr),
    \qquad
    k = 0,1,2,\ldots,
\]
so level $N$ contains $\ell = N, N-2,\ldots,1$ or $0$. For $N = 3$: $\ell = 3,1$, giving $7+3 = 10$ states, matching the Cartesian count.
\end{example}

\begin{fact}[Key results of this chapter]
\leavevmode
\begin{itemize}
\item $E_n = \hbar\omega(n+\frac12)$, Hermite--Gaussian eigenfunctions, zero-point energy.
\item $\hat{a},\hat{a}^\dagger$ with $[\hat{a},\hat{a}^\dagger] = 1$, $\hat{H} = \hbar\omega(\hat{N}+\frac12)$; $\hat{a}^\dagger|n\rangle = \sqrt{n+1}|n+1\rangle$, $\hat{a}|n\rangle = \sqrt n|n-1\rangle$.
\item Moments by counting: $\langle x^2\rangle_n = x_0^2(n+\frac12)$.
\item Thermal oscillators freeze out below $k_BT\sim\hbar\omega$.
\end{itemize}
\end{fact}
